\chapter*{\centering ABSTRAK}
\addcontentsline{toc}{chapter}{ABSTRAK}

\begin{spacing}{1}
\noindent MyPassword adalah perusahaan penyedia solusi teknologi informasi dan integrator sistem, yang menawarkan solusi infrastruktur TI dan aplikasi bisnis. Perusahaan ini terdiri dari beberapa divisi, termasuk \textit{Business Application}, \textit{Infrastructure}, \textit{Human Resources}, dan \textit{Business Support}. Salah satu tantangan yang saat ini dihadapi oleh tim \textit{Human Resources} adalah sistem pengajuan lembur di divisi \textit{Infrastructure}. Saat ini, perusahaan sudah memiliki aplikasi presensi yang memiliki fitur pengajuan lembur. Namun, aplikasi saat ini masih kurang efisien dalam proses pencatatan lembur dan aplikasi tidak bisa dikustomisasi sesuai kebutuhan perusahaan, karena fitur-fiturnya sudah ditentukan oleh pihak eksternal yang mengembangkan aplikasi tersebut. Tim \textit{Human Resources} menginginkan sistem aplikasi yang mampu mencatat dan menghitung lembur secara otomatis ketika \textit{staff} melakukan presensi diluar jam \textit{shift} atau hari libur yang ditetapkan perusahan, baik untuk WFO maupun WFH. Oleh karena itu, dikembangkanlah aplikasi presensi \textit{online} berbasis web dengan menggunakan bahasa PHP dan \textit{framework} Laravel, serta MySQL sebagai basis datanya. Aplikasi presensi \textit{online} berbasis web akan dikembangkan dengan metode Scrum. Aplikasi presensi \textit{online} menggunakan teknik Geolocation untuk mengambil titik koordinat lokasi \textit{staff} ketika melakukan presensi. Titik koordinat yang didapat akan dikonversi menjadi peta interaktif dan nama lengkap jalan. Aplikasi presensi \textit{online} juga mengimplementasikan algoritma Haversine untuk menghitung jarak antara lokasi \textit{staff} saat presensi dengan lokasi kantor MyPassword, hasil perhitungan untuk memberikan keterangan saat presensi. Selain itu, aplikasi juga dilengkapi TensorFlow untuk melakukan validasi wajah saat pengambilan foto presensi. Pengujian fitur aplikasi presensi \textit{online} juga dilakukan oleh \textit{programmer} menggunakan metode \textit{whitebox testing}, sedangkan pengujian oleh \textit{user} menggunakan \textit{blackbox testing}. Pengujian oleh \textit{programmer} akan dilakukan pengujian keseluruhan fitur dari aplikasi termasuk pengujian perhitungan algoritma Haversine. Pengujian perhitungan algoritma Haversine dilakukan dengan membandingkannya terhadap perhitungan \textit{Google Geometry} API pada 20 titik lokasi berbeda. Hasilnya menunjukkan rata-rata perbedaan sebesar 0.0113 km. Selain itu, pengujian oleh \textit{user} akan dilakukan oleh \textit{Staff Infrastructure, Manager Infrastructure}, dan \textit{Manager Human Resources} untuk menguji seluruh fitur sesuai dengan masing-masing \textit{role}. Hasil dari pengujian oleh \textit{programmer} maupun \textit{user} mendapatkan hasil \textit{valid} untuk seluruh \textit{test case}. Pengujian juga dilakukan kepada 10 \textit{staff non-}IT menggunakan metode \textit{blackbox testing} yang mendapatkan hasil valid untuk semua skenario dan pengujian SUS juga dilakukan yang mendapatkan hasil skor 79,5 yang merupakan peringkat A- dalam pengujian SUS.     
\\


 

\noindent\textbf{Kata kunci:} Presensi, lembur, \textit{geolocation}, \textit{framework} Laravel, scrum
\end{spacing}



